% Polarisation measures and their 1-WL status.

\section{Polarisation measures and 1-WL invariance}\label{sec:measures}

This chapter defines each polarisation measure the paper works with
and sorts it into either $\mathcal{F}_{\wl}(\mathcal{G})$ or its
complement. The measures fall into two disjoint groups: measures that
depend on the featured graph only through per-vertex quantities
(degrees, own feature, neighbourhood feature composition) and are
$1$-WL invariant, and measures that depend on the graph through
closure statistics or the Laplacian spectrum and are not.

\subsection{Measures that lie in $\mathcal{F}_{\wl}$}

\begin{definition}[Edge homophily]\label{def:homophily}
Fix a featured graph $(\gph, x)$ with binary or categorical
features. The \emph{edge homophily coefficient} is the fraction of
edges whose endpoints share the same feature:
\[
h(\gph, x) = \frac{|\{(u, v) \in \edges : x_u = x_v\}|}{|\edges|}.
\]
\end{definition}

\begin{definition}[Attribute assortativity]\label{def:assortativity}
Fix a featured graph $(\gph, x)$ with categorical features taking
values in a finite set. The \emph{Newman attribute assortativity
coefficient} is
\[
r(\gph, x) = \frac{\sum_i e_{ii} - \sum_i a_i b_i}{1 - \sum_i a_i
b_i},
\]
where $e_{ij}$ is the fraction of edges from vertices of feature $i$
to vertices of feature $j$, and $a_i, b_i$ are the row and column
sums of the matrix $(e_{ij})$ \citep{Newman2003}.
\end{definition}

\begin{definition}[Modularity of a fixed partition]\label{def:modularity}
Fix a graph $\gph$ and a partition $\mathcal{P}$ of $\vertices$ into
disjoint communities. The \emph{modularity} of $\mathcal{P}$ is
\[
Q(\gph, \mathcal{P}) = \frac{1}{2|\edges|} \sum_{u, v \in \vertices}
\left[ A_{uv} - \frac{\deg(u) \deg(v)}{2|\edges|}\right]
\mathbb{1}[u, v \text{ in the same block of } \mathcal{P}],
\]
where $A_{uv}$ is the adjacency indicator \citep{Newman2004}.
\end{definition}

\begin{proposition}\label{prop:computable-measures}
The edge homophily $h(\gph, x)$, attribute assortativity $r(\gph,
x)$, and modularity $Q(\gph, \mathcal{P})$ of a fixed partition all
lie in $\mathcal{F}_{\wl}(\mathcal{G})$ on every compact featured
graph class $\mathcal{G}$. Any functional of the form
\[
\Psi(\gph, x) = \frac{1}{|\vertices|} \sum_{v \in \vertices} \phi\!\left(x_v,\ \deg(v),\ \{\!\{x_u : u \in \neigh{v}\}\!\}\right)
\]
for continuous $\phi$ also lies in $\mathcal{F}_{\wl}(\mathcal{G})$.
\end{proposition}

\begin{proof}
Edge homophily rewrites as
\[
h(\gph, x) = \frac{1}{2|\edges|} \sum_{v \in \vertices} |\{u \in
\neigh{v} : x_u = x_v\}|,
\]
which is a sum over vertices of a quantity determined by $x_v$ and
the multiset $\{\!\{x_u : u \in \neigh{v}\}\!\}$: exactly the input
to $1$-WL's round-one refinement. Two featured graphs with the same
round-one $1$-WL colour distribution and the same $|\edges|$ give
the same $h$; $|\edges| = \frac{1}{2} \sum_v \deg(v)$ is determined
by the degree multiset, which is a coarser invariant than the
$1$-WL colour multiset. Hence $h$ is $1$-WL invariant.

Attribute assortativity rewrites as a ratio of quadratic forms in
the edge-count matrix $(e_{ij})$, whose entries are determined by
the multiset of $(x_v, \{\!\{x_u : u \in \neigh{v}\}\!\})$ tuples
over $v$. The same argument gives modularity of a fixed partition
(which rewrites through within-partition and cross-partition edge
counts and vertex degrees). The general functional $\Psi$ is
$1$-WL invariant by direct inspection.
\end{proof}

\subsection{Closure-based measures}

Closure statistics count triangles or measure the density of local
neighbourhood links.

\begin{definition}[Local clustering coefficient]\label{def:clustering}
For a vertex $v$ of degree $\deg(v) \ge 2$, the \emph{local
clustering coefficient} is the fraction of pairs of neighbours of
$v$ that are joined by an edge:
\[
C(v) = \frac{2 \cdot |\{\{u, w\} \subseteq \neigh{v} : \{u, w\} \in
\edges\}|}{\deg(v) (\deg(v) - 1)}.
\]
The graph-level \emph{average local clustering} is
\[
\Pi_{\clos}(\gph) = \frac{1}{|\vertices'|} \sum_{v : \deg(v) \ge 2} C(v),
\]
with $\vertices'$ the set of vertices of degree at least $2$.
\end{definition}

\begin{definition}[Triangle density]\label{def:triangles}
The \emph{triangle density} of $\gph$ is the number of triangles in
$\gph$ divided by $\binom{|\vertices|}{3}$:
\[
t(\gph) = \frac{|\{S \in \binom{\vertices}{3} : S \text{ induces a
triangle}\}|}{\binom{|\vertices|}{3}}.
\]
\end{definition}

\subsection{Random-walk controversy and spectral measures}

\begin{definition}[Random walk controversy]\label{def:rwc}
Fix a graph $\gph$ and a bipartition $\{S_+, S_-\}$ of $\vertices$
into two disjoint communities. Fix a small integer $k$ and identify
\emph{authorities}: the $k$-highest-degree vertices of each side.
For $\sigma, \tau \in \{+, -\}$, define $P_{\sigma\tau}$ as the
probability that a simple random walk starting at a uniform vertex
of $S_\sigma$ hits an authority of $S_\tau$ before any other
authority. The \emph{random walk controversy} of $(\gph, \{S_+,
S_-\})$ is
\[
\mathrm{RWC}(\gph, \{S_+, S_-\}) = P_{++} P_{--} - P_{+-} P_{-+}
\in [-1, 1]
\]
\citep{Garimella18}.
\end{definition}

\begin{definition}[Generalised Euclidean distance]\label{def:ged}
Fix a graph $\gph$ and an opinion vector $o \in [-1, +1]^{|\vertices|}$.
Split $o$ into a non-negative part $o^+_v = \max(o_v, 0)$ and a
non-positive part $o^-_v = \max(-o_v, 0)$. Let $L$ be the
combinatorial Laplacian of $\gph$ and $L^\dagger$ its Moore-Penrose
pseudoinverse. The \emph{generalised Euclidean distance} of
\citet{HohmannGED2023} is
\[
\mathrm{GED}(\gph, o) = \sqrt{(o^+ - o^-)^\top L^\dagger (o^+ - o^-)}.
\]
On a disconnected graph, $L^\dagger$ is applied to each connected
component separately; \citet{HohmannGED2023} restrict to connected
graphs in their formulation.
\end{definition}

\subsection{Three closure and spectral measures are not $1$-WL invariant}

\begin{proposition}\label{prop:noncomputable-measures}
The functionals $\Pi_{\clos}$, $t$, and $\mathrm{GED}$ under a
non-trivial opinion vector each fail $1$-WL invariance on any class
$\mathcal{G}$ containing the pair $(A, B) = (C_3 \sqcup C_3, C_6)$
of Section~\ref{sec:warmup}. In particular, none of the three lies in
$\mathcal{F}_{\mpnn}(\mathcal{G})$ on such a class.
\end{proposition}

\begin{proof}
Recall $A \sim_{\wl} B$ from Section~\ref{sec:warmup} in the unfeatured
case. In the featured versions with a constant feature $\bot$, the
same argument gives $(A, \bot) \sim_{\wl} (B, \bot)$.

\emph{Parts (a), (b): clustering and triangle density.} $\Pi_{\clos}(A)
= 1$ and $\Pi_{\clos}(B) = 0$ from Section~\ref{sec:warmup}; the triangle
counts are $2$ and $0$, so $t(A) = 2/\binom{6}{3} = 1/10$ and $t(B)
= 0$. The two graphs are $1$-WL equivalent and the functionals
differ on them.

\emph{Part (c): GED.} Take the opinion vector $o = (+1, +1, +1, -1,
-1, -1)^\top$. The computation in Appendix~\ref{app:ged-witness} gives
$\mathrm{GED}(A, o) = 0$ and $\mathrm{GED}(B, o) = \sqrt{11/18}$
using the per-component convention of \citet{HohmannGED2023}. The
values differ.

By Proposition~\ref{prop:upper}, none of the three lies in
$\mathcal{F}_{\mpnn}(\mathcal{G})$ on a class containing the
witness $(A, B)$.
\end{proof}

\begin{remark}[Status of the random walk controversy score]
\label{rem:rwc}
The random walk controversy score of \citet{Garimella18} is a
functional of a graph together with a bipartition, not of the
graph alone. Whether $\mathrm{RWC}$ lies in
$\mathcal{F}_{\wl}(\mathcal{G})$ depends on how the bipartition is
supplied. If the bipartition is encoded as a $\{-1, +1\}$
featured-graph attribute, the featured version of the natural
witness pair $(C_3 \sqcup C_3, C_6)$ under matching bipartitions is
not $1$-WL equivalent, so this pair does not witness non-invariance
of $\mathrm{RWC}$. We attempted the alternative witness of the
Shrikhande and $4 \times 4$ rook's graphs (both strongly regular with
parameters $(16, 6, 2, 2)$, hence $1$-WL equivalent
\citep{Shrikhande1959}). Under a bipartition into rows of size $8$
that preserves featured $1$-WL equivalence, exact absorbing-chain
computation gives $\mathrm{RWC}(R) = \mathrm{RWC}(S) = 0.0476$;
the two values coincide. The apparent reason is that both graphs
are vertex-transitive and cospectral, and the symmetry propagates
to hitting probabilities under symmetric bipartitions. A rigorous
witness for $\mathrm{RWC} \notin \mathcal{F}_{\mpnn}(\mathcal{G})$
requires a pair of graphs, a bipartition-as-feature under which the
pair is $1$-WL equivalent, and a computation on which $\mathrm{RWC}$
differs. Constructing such a witness is left as an open technical
problem; this paper does not treat $\mathrm{RWC}$'s $1$-WL
invariance status as established.
\end{remark}

\subsection{The taxonomy}

The consequence of Theorem~\ref{thm:dichotomy} for the polarisation
measurement literature is a two-column taxonomy of the published
measures.

\begin{center}
\begin{tabular}{ll}
\hline
$\mathcal{F}_{\wl}$-invariant (MPNN-computable) & Not $1$-WL invariant \\
\hline
Edge homophily $h$ & Local clustering $\Pi_{\clos}$ \\
Attribute assortativity $r$ & Triangle density $t$ \\
Modularity $Q$ (fixed partition) & Generalised Euclidean distance $\mathrm{GED}$ \\
Any $\Psi$ from Proposition~\ref{prop:computable-measures} & \\
\hline
\multicolumn{2}{l}{Random walk controversy $\mathrm{RWC}$: status
open; see Remark~\ref{rem:rwc}.} \\
\hline
\end{tabular}
\end{center}

A message-passing GNN estimator can approximate every functional in
the left column arbitrarily well on the appropriate compact
featured-graph class. No message-passing GNN can approximate any
functional in the right column on any class containing the six-vertex
pair $(A, B)$. The random walk controversy score plausibly belongs on
the right, but this paper does not construct a witness. The next
chapter shows the impossibility survives on random-graph ensembles
of arbitrary size.
